\SourceLeaf{369}%
\noindent be the cosines of the three angles that a straight line $0P$ makes with $01,02,03$, respectively. This straight line has the same position in the moving space after the mutation as before it. [If we further put]
\[
\BareRoot{(aa+\rho\rho)}=k
\]
or
\[
k=\BareRoot{(aa+bb+cc+dd)},
\]
then $k$ is the linear enlargement of the mutated space.

Finally, on putting
\[
a=k\cos\theta,\qquad \rho=k\sin\theta
\]
we find that $2\theta$ is the rotation about $0P$ which, combined with the enlargement $k$, can produce the mutation.

\Needspace{12\baselineskip}
\begin{center}\bfseries REMARK.\end{center}

The notes grouped under [I] are in a notebook. Numbers 1 to 7 agree exactly in handwriting and ink with astronomical records from 1819 on the preceding pages. The remark designated number 8 looks different. In substance it has the character of a later addition. This supposition is also confirmed by the fact that astronomical notes apparently written at the same time as it date from 1822 to 1823. Note [II] fills a single undated slip. Judging by the handwriting, it dates from a still later time.
\begin{flushright}\textsc{Stäckel.}\end{flushright}
